\section{Quantitative Gaussian edge estimates}\label{sec:quantitative-edge}

The following estimates quantify the two scales without a relation between
their two small errors.  The proof uses the mean under each displaced
Gaussian measure.  In particular, it does not divide two estimates whose
errors are only small on the common-amplitude scale.

\begin{theorem}[Uniform errors for two Chebyshev families]
\label{thm:quantitative-edge}
Fix $0<s_0\le s_1<\infty$ and $0\le C_0<\infty$.  Consider either
\[
 p_n=T_n\quad(n\ge1),\qquad\text{or}\qquad
 p_n=U_{n-1}/n\quad(n\ge1),
\]
where $U_{n-1}(\cos\theta)=\sin(n\theta)/\sin\theta$.
Put
\[
 Q^p_{n,A}(s,c)=(e^{sA^{-2}D^2}p_n)(1+c/A),\quad
 \mu=(n^2/A)^{1/3},\quad \epsilon=\mu/A,
\]
and set $F_s=3\,2^{-4/3}s^{1/3}$ and
$G_s=2^{-2/3}s^{-1/3}$.
There are constants $C,\mu_*>0$ and $\epsilon_*>0$, computable from
$s_0,s_1,C_0$ alone, such that if $\mu\ge\mu_*$ and
$0<\epsilon\le\epsilon_*$, then $Q^p_{n,A}(s,c)>0$ and
\begin{align}
 \left|\mu^{-2}\log Q^p_{n,A}(s,0)-F_s\right|
 &\le C\left(\epsilon+
             \frac{\log(2+\mu)}{\mu^2}\right),
 \label{eq:quantitative-amplitude}\\
 \left|\mu^{-1}\log\frac{Q^p_{n,A}(s,c_2)}{Q^p_{n,A}(s,c_1)}
                -G_s(c_2-c_1)\right|
 &\le C|c_2-c_1|\left(\sqrt\epsilon+
             \frac{\sqrt{\log(2+\mu)}}\mu\right).
 \label{eq:quantitative-tilt}
\end{align}
These estimates hold uniformly for $s\in[s_0,s_1]$ and
$|c_1|,|c_2|\le C_0$.  No optimality of these error orders is asserted.
In particular, they apply throughout $A\to\infty$, $\mu\to\infty$,
$\mu/A\to0$.
\end{theorem}

For auditability, a profile proposition gives explicit constants before
they are absorbed into $C$.

\begin{proposition}[A square-root edge profile criterion]
\label{prop:edge-profile}
Write $\psi(z)=\sqrt{2z}$, $z_-=2^{1/3}s_0^{2/3}$,
$z_+=2^{1/3}s_1^{2/3}$, $Z=z_++1$, and $J=[z_-,Z]$.
Suppose a real measurable profile $p=p_{\mu,A}$, with $A=\mu/\epsilon$,
satisfies the following inequalities for some fixed $k>0$, $b,d\ge0$:
\begin{align}
 |p(x)|&\le1 &&(|x|\le1),\label{eq:profile-interior}\\
 0<p(1+\mu z/A)&\le e^{\mu^2\psi(z)},\quad
 |p(-1-\mu z/A)|\le e^{\mu^2\psi(z)}
 &&(z\ge0),\label{eq:profile-envelope}\\
 p(1+\mu z/A)&\ge k\mu^{-b}
       e^{\mu^2(\psi(z)-d\epsilon)} &&(z\in J).
 \label{eq:profile-local-lower}
\end{align}
Assume these hold for $\mu\ge\mu_0\ge1$ and
$0<\epsilon\le\epsilon_0\le1$.  Define $Q^p$ by its Gaussian
convolution, which agrees with the finite heat expansion for a polynomial:
\[
 Q^p(s,c)=\frac1{\sqrt{4\pi s}}\int_{\mathbb R}
       p(1+y/A)e^{-(y-c)^2/(4s)}\,dy.
\]
The conclusions of Theorem~\ref{thm:quantitative-edge} hold.
Constants may also depend on $k,b,d,\mu_0,\epsilon_0$.
The formulas below specify the constants and a noncircular threshold.
\end{proposition}

\paragraph{Constants for the profile criterion.}
In the following formulas, $\log_+x=\max\{0,\log x\}$.  All constants
are fixed before $\mu,A$ are chosen.  Set
\begin{align*}
 F_*&=3\,2^{-4/3}s_0^{1/3},&
 D_*&=(2z_-)^{-3/2}+(2s_0)^{-1},& a_*&=(8s_1)^{-1},\\
 H_*&=C_0^2s_1/(2s_0^2),&
 L_*&=\frac{k}{\sqrt{4\pi s_1}}
              e^{-D_*/2-C_0/(2s_0)},&
 V_*&=\sqrt{2s_1/s_0}\,e^{H_*},\\
 K_*&=\frac{e^{H_*}}{\sqrt{4\pi s_0}L_*},&
 H_{\rm tail}&=3\,2^{-4/3}(2s_1)^{1/3},&
 E_0&=1+\sqrt{s_1/(2s_0)},\\
 E_1&=\frac1{\sqrt{\pi s_0}}+
 \frac{6\sqrt{s_1/e}\sqrt{2\pi s_1}+4s_1}
      {2s_0\sqrt{4\pi s_0}},&
 \kappa&=F_*/4.&&
\end{align*}
For $\mu\ge1$ put
\begin{align*}
 r(\mu,\epsilon)^2
  &=\frac{d\epsilon}{a_*}
    +\frac{(b+1)\log\mu+\log_+(K_*/a_*)}{a_*\mu^2},\\
 \rho(\mu,\epsilon)&=r(\mu,\epsilon)+\mu^{-2},\\
 \rho_*&=\sqrt{\frac{d+b+1+\log_+(K_*/a_*)}{a_*}}+1,\\
 M_*&=(z_++\rho_*+C_0)/(2s_0),\qquad
 E_2=2(E_1+E_0M_*).
\end{align*}
One admissible threshold is
\begin{align}
 \epsilon_*&=\min\left\{\epsilon_0,\,
             \frac1{\sqrt{8s_1H_{\rm tail}}},\,
             \frac{F_*}{4d}\right\},\label{eq:profile-epsilon-threshold}\\
 \mu_*&=\max\left\{2,C_0,\mu_0,
       \frac{2C_0z_+}{s_0F_*},\frac{8b}{F_*},
       \sqrt{\frac{8\{\log_+(1/L_*)+C_0^2/(4s_0)\}}{F_*}},
       \sqrt{\frac{\log(2E_0)}\kappa}\right\}.
       \label{eq:profile-mu-threshold}
\end{align}
The entry $F_*/(4d)$ is omitted when $d=0$.
At this threshold, more explicit versions of the two errors are
\begin{align}
 \left|\mu^{-2}\log Q^p(s,0)-F_s\right|
 &\le d\epsilon+
 \frac{b\log\mu+\max\{\log_+(1/L_*),\log_+V_*\}
              +2E_0e^{-\kappa\mu^2}}{\mu^2},
 \label{eq:profile-amplitude-explicit}\\
 \left|\mu^{-1}\log\frac{Q^p(s,c_2)}{Q^p(s,c_1)}
                 -G_s(c_2-c_1)\right|
 &\le |c_2-c_1|\left\{
       \frac{\rho(\mu,\epsilon)+C_0/\mu}{2s_0}
                     +E_2e^{-\kappa\mu^2}\right\}.
 \label{eq:profile-tilt-explicit}
\end{align}

\begin{proof}[Proof of Proposition~\ref{prop:edge-profile}]
Retain the positive exterior and write
\[
 I(s,c)=\frac1{\sqrt{4\pi s}}\int_0^\infty
       p(1+y/A)e^{-(y-c)^2/(4s)}\,dy.
\]
Let $z_s=2^{1/3}s^{2/3}$ and
$\Phi_s(z)=\psi(z)-z^2/(4s)$.  Its derivative vanishes at $z_s$,
$\Phi_s(z_s)=F_s$, and
\[
 \Phi_s''(z)=-(2z)^{-3/2}-(2s)^{-1}\le-(2s_1)^{-1}
 \quad(z>0).
\]
Consequently, including $z=0$ by continuity,
\begin{equation}\label{eq:profile-strong-concavity}
 \Phi_s(z)\le F_s-\frac{(z-z_s)^2}{4s_1}\quad(z\ge0).
\end{equation}
Put $\mathcal B=\mu^2F_s+c\mu z_s/(2s)-c^2/(4s)$.
On $z_s\le z\le z_s+1/\mu$, the interval lies in $J$,
$|\Phi_s''(z)|\le D_*$, and the linear tilt changes by at most
$C_0/(2s_0)$.  Integrating the local lower bound over this interval
therefore gives
\begin{equation}\label{eq:profile-I-lower}
 I(s,c)\ge L_*\mu^{-b}e^{\mathcal B-d\mu^2\epsilon}.
\end{equation}
For the upper bound, Young's inequality in the precise form
\[
 \frac{c\mu(z-z_s)}{2s}
 \le\frac{\mu^2(z-z_s)^2}{8s_1}+H_*
\]
and \eqref{eq:profile-strong-concavity} show that
\begin{equation}\label{eq:profile-I-upper}
 I(s,c)\le V_*e^{\mathcal B}.
\end{equation}
Indeed the changed measure is $\mu\,dz/\sqrt{4\pi s}$ and its Gaussian
upper bound can be integrated over the whole real line.

Normalize the positive integrand for each $c$ to a probability measure
for $Z=y/\mu$.  Equations~\eqref{eq:profile-I-lower} and
\eqref{eq:profile-strong-concavity} give its density bound
\[
 f_{s,c}(z)\le K_*\mu^{b+1}e^{d\mu^2\epsilon}
                  e^{-a_*\mu^2(z-z_s)^2}\quad(z\ge0).
\]
For any $r\ge0$, extending both tails to the whole line gives
\[
 \mathbb E_{s,c}|Z-z_s|
 \le r+\frac{K_*}{a_*}\mu^{b-1}
          e^{d\mu^2\epsilon-a_*\mu^2r^2}.
\]
The displayed choice of $r(\mu,\epsilon)$ makes the second term at
most $\mu^{-2}$.  Differentiating the Gaussian, justified by the
global envelope, now yields the exact identity
\[
 \partial_c\log I(s,c)=\frac{\mu\mathbb E_{s,c}Z-c}{2s}.
\]
It follows that
\begin{equation}\label{eq:profile-positive-derivative}
 \left|\mu^{-1}\partial_c\log I(s,c)-G_s\right|
       \le\frac{\rho(\mu,\epsilon)+C_0/\mu}{2s_0}.
\end{equation}
This estimate treats the unbounded Gaussian tilt directly and uniformly
for positive and negative $c$.

It remains to restore the signed remainder $\mathcal R=Q^p-I$.
Its interior contribution from $-2A\le y\le0$ has absolute value at
most one and $c$-derivative at most $(\pi s_0)^{-1/2}$.  In the
negative exterior put $y=-2A-u$, $u\ge0$.  Since $A\ge C_0$,
\[
 2A+u+c\ge A+u,\qquad
 \mu^2\psi(u/\mu)-\frac{(2A+u+c)^2}{4s}
 \le-\frac{A^2}{4s_1}+H_{\rm tail}\mu^2
                            -\frac{u^2}{8s_1}.
\]
The value of $H_{\rm tail}$ follows by maximizing
$\sqrt{2z}-z^2/(8s_1)$ on $z\ge0$; its maximizer is
$2^{1/3}(2s_1)^{2/3}$.  The threshold for $\epsilon$ makes the
right side at most $-A^2/(8s_1)-u^2/(8s_1)$.
Gaussian integration proves $|\mathcal R|\le E_0$.
For its derivative use
$|2A+u+c|/(2s)\le(3A+u)/(2s_0)$, together with
\[
 \int_0^\infty e^{-u^2/(8s_1)}du=\sqrt{2\pi s_1},\quad
 \int_0^\infty u e^{-u^2/(8s_1)}du=4s_1,\quad
 \sup_{A\ge0}Ae^{-A^2/(8s_1)}=2\sqrt{s_1/e}.
\]
This proves $|\partial_c\mathcal R|\le E_1$ with the stated constant.
No assumption on the sign of either omitted portion is used.

The first threshold gives $d\epsilon\le F_*/4$.  The linear term
in $\mathcal B$ costs at most $F_*\mu^2/4$ by the fourth entry of
\eqref{eq:profile-mu-threshold}.  Finally $\log\mu\le\mu$,
the fifth and sixth entries, and \eqref{eq:profile-I-lower} give
\[
 I(s,c)\ge e^{\kappa\mu^2},\qquad
 |\mathcal R/I|\le E_0e^{-\kappa\mu^2}\le\tfrac12.
\]
Thus $Q^p>0$ and
$|\log Q^p-\log I|\le2E_0e^{-\kappa\mu^2}$.
Since $\epsilon\le1$ and $\log\mu\le\mu^2$, we have
$\rho\le\rho_*$ and $|\partial_c\log I|\le\mu M_*$.  Hence
\[
 \mu^{-1}|\partial_c\log Q^p-\partial_c\log I|
 =\mu^{-1}\left|
   \frac{\partial_c\mathcal R-\mathcal R\partial_c\log I}
        {I+\mathcal R}\right|
 \le E_2e^{-\kappa\mu^2}.
\]
Integrating this estimate and
\eqref{eq:profile-positive-derivative} between $c_1$ and $c_2$
proves \eqref{eq:profile-tilt-explicit}, including exact vanishing when
$c_1=c_2$.  At $c=0$, use \eqref{eq:profile-I-lower} and
\eqref{eq:profile-I-upper} to obtain
\eqref{eq:profile-amplitude-explicit}.  These explicit estimates imply
the simpler displayed orders in the theorem.  For example,
$e^{-\kappa\mu^2}\le(2e\kappa)^{-1/2}/\mu$ for $\mu>0$,
and each remaining constant term is absorbed using
$\log(2+\mu)\ge1$ for $\mu\ge2$.
\end{proof}

\begin{proof}[Verification of Theorem~\ref{thm:quantitative-edge}]
Write $\theta=\operatorname{arcosh}(1+\epsilon z)$.  For $u\ge0$,
integrating $(2t)^{-1/2}(1+t/2)^{-1/2}$ from zero to $u$ gives
\begin{equation}\label{eq:arcosh-quantitative-bounds}
 \sqrt{2u}(1-u/12)\le\operatorname{arcosh}(1+u)
 \le\sqrt{2u},\qquad
 \operatorname{arcosh}(1+u)\ge
                \frac{\sqrt{2u}}{\sqrt{1+u/2}}.
\end{equation}
The first lower bound uses $(1+v)^{-1/2}\ge1-v/2$; the other uses
the minimum of the decreasing factor in the same integral.
As $n\sqrt\epsilon=\mu^2$, it follows on $J$ that
\[
 \mu^2\{\psi(z)-d\epsilon\}\le n\theta
       \le\mu^2\psi(z),\qquad d=\sqrt2 Z^{3/2}/12.
\]
For $T_n$, the identity $T_n(\cosh\theta)=\cosh(n\theta)$ gives
the exterior envelope and local lower bound with
$k=1/2$, $b=0$, $\mu_0=1$, $\epsilon_0=1$.
Its interior bound and parity give the other two hypotheses.

For the second family, the finite sum
\[
 \frac{U_{n-1}(\cosh\theta)}n
  =\frac1n\sum_{j=0}^{n-1}e^{(n-1-2j)\theta}
  =\frac{\sinh(n\theta)}{n\sinh\theta}
\]
is positive and at most $e^{n\theta}$.  The corresponding trigonometric
sum has modulus at most one, proving the interior bound, including the
endpoints by continuity.  Parity gives the negative exterior envelope.
On $J$ one has the exact denominator identity
\[
 n\sinh\theta=\mu^2\sqrt{z(2+\epsilon z)}
       \le\mu^2\sqrt{Z(2+Z)}.
\]
Put $h_*=\sqrt{2z_-}/\sqrt{1+Z/2}$ and
$\mu_0=\max\{1,h_*^{-1/2}\}$.  Then
$n\theta\ge\mu^2h_*\ge1$ on $J$ by
\eqref{eq:arcosh-quantitative-bounds}.  Therefore
\[
 \frac{U_{n-1}(1+\epsilon z)}n
 \ge\frac{1-e^{-2}}{2\sqrt{Z(2+Z)}}\mu^{-2}
                    e^{\mu^2\{\psi(z)-d\epsilon\}}.
\]
This verifies the criterion with $b=2$,
$k=(1-e^{-2})/(2\sqrt{Z(2+Z)})$, and $\epsilon_0=1$.
All constants and thresholds are consequently obtained from the explicit
formulas above.  Taking their maximum or minimum as appropriate gives
one choice valid for both families.
\end{proof}

\begin{remark}[Normalization and the scope of the extension]
For unnormalized $U_{n-1}$ the displacement ratio is unchanged, but
the common-amplitude formula contains the additional term
$\mu^{-2}\log n$.  This term need not vanish in the original joint
range.  For example, take $n_j=\lceil e^{j^3}\rceil$ and
$A_j=n_j^2/j^3$.  Then $\mu_j=j$, $\epsilon_j=j^4/n_j^2\to0$,
whereas $\mu_j^{-2}\log n_j\sim j$.
The normalization therefore removes a real obstruction, rather than
only changing notation.  The profile criterion requires positivity on
the exterior, a global envelope on both sides, an interior bound, and
a quantitative lower bound near the saddle.  It does not assert the
same conclusion for arbitrary Jacobi parameters or arbitrary
polynomial heat flows.  The normalized second family has the same two
leading constants; no Navier--Stokes realization for that additional
family is claimed here.
\end{remark}
